\section{Related Work}

The literature on automated crowd anomaly detection spans more than a decade of active research. In this section, we briefly overview the related literature, which includes methods that leverage optical flow statistics, deep temporal models, and explainable AI surveillance systems. The most relevant prior works are summarized in Table~\ref{tab:literature}.

\subsection{Optical-Flow-Based Motion Analysis}

The most essential part of crowd behavior analysis is motion activity. Optical flow allows to estimate motion patterns between consecutive frames by calculating displacement vectors for each pixel. This information can be utilized to build descriptors that summarize the motion field statistics, such as distribution of magnitude and direction of displacement vectors. Optical flow is advantageous over other features in that it is independent of the appearance of objects in a video.

\subsubsection{Dense Optical Flow Algorithms}
The choice of a particular algorithm has little effect on the final results as most methods utilize similar motion estimation statistics. Sparse algorithms such as Lucas-Kanade track only salient points and have fewer computational costs, but fail to estimate motion in areas with little texture; dense algorithms calculate displacement vectors for all frames.

The Farneb\"{a}ck algorithm \cite{farneback2003two}, which is employed throughout the present work, is a dense optical flow method that approximates the local neighbourhood of each pixel as a quadratic polynomial and solves for the displacement that best explains the polynomial coefficient changes between frames. It offers a favourable trade-off between accuracy and computational cost: it produces reliable flow estimates in the texture-poor backgrounds common in surveillance footage (e.g., plain walls, uniform floors) where sparse methods fail, while remaining fast enough for real-time processing. Alternative dense methods include Horn-Schunck (which imposes a global smoothness constraint but is sensitive to large displacements) and deep-learning-based estimators such as FlowNet and RAFT (which achieve higher accuracy but at substantially greater computational cost, making them unsuitable for edge deployment).

\subsubsection{Feature Engineering from Flow Fields}
Early attempts at optical-flow feature engineering began with Mehran et al.\ \cite{mehran2009abnormal}, who computed attractive and repulsive forces from flow fields to model social force equations. However, this approach assumes relatively sparse crowds where individual force vectors can be estimated, limiting its use in high-density settings. 

Subsequently, the approach of extracting handcrafted scalar statistics from dense optical flow fields was formalised for stampede detection by Cob-Parro et al.\ \cite{cobparro2021stampede}. Their pipeline computes three descriptors per temporal window:
\begin{itemize}
  \item \textbf{Entropy} of the flow-direction histogram, which quantifies the disorder of movement directions. In normal crowd flow, entropy is low because most people move in a shared direction. During the period of panic breakouts, entropy of motion is high due to people moving chaotically.
  \item \textbf{Temporal Occupancy Variation (TOV)}: TOV is the standard deviation of the occupancy mask thresholded at a certain percentile of flow magnitude. This statistic captures the spatial spreading of an abnormal flow: when a crowd is moving irregularly, the area occupied by motion grows.
  \item \textbf{Kernel Density Estimation (KDE)}: KDE is the kernel density estimation of the flow magnitude distribution, which is a non-parametric way of estimating the probability density function. The distribution is shifted towards higher values during stampede events, as the crowd movements involve higher speeds.
\end{itemize}

These three descriptors are useful in classifying panic events, but they lack a statistic that utilizes the directional information of the motion. Consider a case when a group of pedestrians walks quickly towards a specific location with little deviation from the mean direction. Such a scenario is similar to a stampede in that the crowd moves purposefully, but it does not present any danger. In this case, the computed entropy would be low, TOV would show only a small increase, and KDE would be shifted to higher values. Due to the similarity of these statistics to the ones of panic events, misclassification would be likely. The directional spread of such a regular flow would have lower values than that of a stampede. None of the available descriptors explicitly considers the directional spread, which limits the performance of prior methods.

Chen et al.\ \cite{chen2023opticalflow} experimented with higher order optical flow statistics and concluded that the spread of the magnitude distribution contributes significantly to the ability to detect anomalies. Particularly, the standard deviation of the magnitude is useful in characterizing irregular crowd motion. In addition, they observed that the skewness of the distribution demonstrates the tendency of the motion magnitude. Based on these findings, we incorporated the mean, standard deviation, and skewness of the flow magnitude into our set of descriptors.

\subsection{Deep Temporal Models for Sequence Classification}

When analyzing crowd motion, it is important to consider that an unusual event usually takes some time to unfold. That is, a single frame rarely contains enough information for accurate classification. Effective classification requires modelling the evolution of crowd-state features over time.

\subsubsection{Long Short-Term Memory (LSTM) Networks}
The LSTM architecture \cite{hochreiter1997long} addressed the vanishing gradient problem that limited the ability of vanilla recurrent neural networks to capture long-range temporal dependencies. It introduced a memory cell regulated by three gating mechanisms---input, forget, and output gates---that control the flow of information into, within, and out of the cell. This design allows LSTMs to selectively retain information over extended sequences, making them well-suited for processing the temporal feature vectors generated by optical-flow analysis. The Gated Recurrent Unit (GRU) \cite{cho2014gru} later simplified the LSTM design by merging the forget and input gates into a single ``update'' gate, reducing the parameter count with minimal impact on empirical performance.

\subsubsection{Bidirectional Processing}
A standard LSTM processes the input sequence from past to future, meaning that the hidden state at any time-step encodes only information from preceding frames. Bidirectional variants \cite{chung2014bilstm} address this limitation by running a second LSTM in the reverse direction and concatenating the forward and backward hidden states at each time-step. The resulting representation captures context from both temporal directions, which is particularly valuable for stampede detection: the few frames immediately before panic onset appear ambiguous when viewed in isolation, but become clearly recognisable when the model has access to information about the escalation that follows.

\subsubsection{Attention Mechanisms}
The attention mechanism, originally proposed by Bahdanau et al.\ \cite{bahdanau2015neural} for neural machine translation, provides a principled method for assigning non-uniform importance to different elements of a sequence. Rather than compressing the entire sequence into a single fixed-length context vector (which forces all time-steps to contribute equally), attention computes a weighted sum over all hidden states, where the weights are learned through a small feed-forward network. Time-steps that carry more discriminative information receive higher weights.

Wang et al.\ \cite{wang2023anomaly} represented crowds as spatio-temporal graphs and applied graph convolutional networks with temporal self-attention. While effective for capturing inter-person dynamics, graph-based methods depend on reliable individual detection, which degrades in dense crowd scenarios due to occlusion. Yuan et al.\ \cite{yuan2024crba} combined DenseNet201 features with a BiLSTM and multi-head attention. Their system achieved competitive results on standard benchmarks. However, the reliance on DenseNet201 as a spatial feature extractor introduces approximately 20 million parameters, resulting in high computational cost and limiting deployment on resource-constrained hardware. Li et al.\ \cite{li2024transformer} explored a fully Transformer-based architecture for crowd-flow prediction, replacing recurrent layers entirely with global self-attention. While Transformers capture long-range dependencies more effectively than LSTMs, their quadratic memory complexity with respect to sequence length and higher parameter counts make them impractical for real-time edge deployment.

\subsection{Explainable AI in Surveillance Systems}

The deployment of automated decision-support systems in safety-critical surveillance settings raises important questions about model interpretability and operator trust. A system that flags a crowd event as ``anomalous'' without providing any justification for the decision is difficult to integrate into operational workflows.

Lundberg and Lee \cite{lundberg2017unified} introduced SHAP (SHapley Additive exPlanations), a unified framework for interpreting model predictions based on Shapley values. In the context of crowd anomaly detection, a related technique---permutation feature importance---is employed in the present work to identify which optical-flow features contribute most to classification performance.

Attention-weight visualisation provides a complementary form of explainability at the temporal level. The learned attention weights $[\alpha_1, \ldots, \alpha_S]$ directly indicate which time-steps within the input sequence the model considers most informative for its classification decision, answering the question ``\textit{when} did the model detect the onset of abnormal behaviour?''---information that is critical for deciding how urgently to respond.

\subsection{Summary and Positioning of the Present Work}

The review presented above reveals a clear gap at the intersection of three research threads. Generative models like MDT \cite{mahadevan2010anomaly} or reconstruction methods like sparse coding \cite{lu2013abnormal} have been explored, but they remain sensitive to spatial layout. Optical-flow-based methods offer scene-independent, lightweight features but rely on insufficient descriptor sets that lack measures of directional coordination. Deep temporal models (BiLSTM, attention) have been applied but predominantly with heavyweight spatial feature extractors that preclude real-time edge deployment.

As noted in the comprehensive survey by Liu et al.\ \cite{liu2024survey}, real-time latency and cross-dataset generalisation remain two major bottlenecks in the literature. In this work, we combine the comprehensive set of 9 optical flow statistics described above with a BiLSTM+Attention classifier and an interpretability module.

\begin{table*}[!t]
\centering
\caption{Comparative summary of crowd anomaly detection methods. The proposed approach is the only method to simultaneously address explainability, early detection, robustness evaluation, and cross-dataset generalisation.}
\label{tab:literature}
\footnotesize
\begin{tabular}{p{4.2cm}p{3.2cm}p{2.4cm}ccc}
\toprule
\textbf{Paper} & \textbf{Features} & \textbf{Model} & \textbf{Explain.} & \textbf{Early Det.} & \textbf{Robust.} \\
\midrule
Mehran et al.~\cite{mehran2009abnormal}      & Social Force + OF  & SVM        & No  & No  & No \\
Mahadevan et al.~\cite{mahadevan2010anomaly} & MDT                & Generative & No  & No  & No \\
Lu et al.~\cite{lu2013abnormal}              & Sparse Codes       & Classifier & No  & No  & No \\
Cob-Parro et al.~\cite{cobparro2021stampede} & Entropy+TOV+KDE    & LSTM       & No  & No  & No \\
Wang et al.~\cite{wang2023anomaly}           & GCN + Self-Attn    & Graph NN   & Partial & No & No \\
Yuan et al.~\cite{yuan2024crba}              & DenseNet Features  & BiLSTM+Attn & No & No & No \\
Chen et al.~\cite{chen2023opticalflow}       & OF Statistics      & CNN-LSTM   & No  & No  & No \\
Li et al.~\cite{li2024transformer}           & Flow + Spatial     & Transformer & No & No & No \\
\textbf{Proposed}                            & \textbf{9 OF Stats} & \textbf{BiLSTM+Attn} & \textbf{Yes} & \textbf{Yes} & \textbf{Yes} \\
\bottomrule
\end{tabular}
\end{table*}
